%%%
%
% $Autor: Wings $
% $Datum: 2021-05-14 $
% $Pfad:  $
% $Dateiname: SBOM
% $Version: 4620 $
%
% !TeX spellcheck = de_DE/GB
% !TeX program = pdflatex
% !BIB program = biber/bibtex
% !TeX encoding = utf8
%
%%%




\chapter{SBOM}

%------------------------------------------------

% SOFTWARE SECTION

%------------------------------------------------

\subsection{Software Bill of Materials}

The software environment used in this project supports data preprocessing, time-series modelling, visualization, and reproducible experimentation. The tools listed below form the complete development environment used for implementing the CPI forecasting models.

\begin{table}[h]
	\centering
	\begin{tabular}{|p{3cm}|p{7cm}|p{2cm}|p{2cm}|}
		\hline
		\textbf{Software} & \textbf{Description} & \textbf{Quantity} & \textbf{Price} \\ \hline
		Visual Studio Code & Lightweight and extensible source code editor used for Python development, debugging, and project implementation. & 1 & Free \\ \hline
		Python & Primary programming language used for data preprocessing, statistical modelling, and CPI forecasting implementation. & 1 & Free \\ \hline
		Git & Distributed version control system used to track code changes and manage project versions. & 1 & Free \\ \hline
		GitHub & Online platform used to host the project repository and maintain project documentation and version history. & 1 & Free \\ \hline
		Python Libraries & Libraries including Pandas, NumPy, Matplotlib, Statsmodels, and Scikit-learn used for data processing, visualization, and forecasting models. & 1 & Free \\ \hline
	\end{tabular}
	\caption{Software Bill of Materials}
\end{table}

\subsection{Python Environment}

Python is the primary programming language used in this project. It provides a large ecosystem of libraries for statistical modelling, machine learning, and data analysis.

\begin{itemize}
	\item Python Version: Python 3.11 or later
	\item Package Manager: pip
	\item Libraries Used:
	\begin{itemize}
		\item Pandas – Data manipulation and time-series handling
		\item NumPy – Numerical computation
		\item Matplotlib – Data visualization
		\item Statsmodels – Implementation of ARIMA, SARIMA, and ETS forecasting models
		\item Scikit-learn – Evaluation metrics and model validation
	\end{itemize}
\end{itemize}

\subsection{Visual Studio Code (VS Code)}

Visual Studio Code (VS Code) is used as the primary development environment for implementing the CPI forecasting models.

\begin{enumerate}
	
	\item User Interface: VS Code provides a customizable interface that supports themes, layouts, and workspace configuration.
	
	\item Extensions and Marketplace: Extensions such as the Python extension and Git integration provide additional development tools and debugging support.
	
	\item Integrated Terminal: Developers can execute Python scripts and command-line operations directly within the editor.
	
	\item IntelliSense: Provides intelligent code completion, parameter hints, and syntax suggestions.
	
	\item Debugging Support: Enables breakpoints, variable inspection, and step-by-step execution of Python programs.
	
	\item Version Control Integration: Built-in Git integration allows efficient version control and code management.
	
\end{enumerate}

\subsection{Version Control and Documentation}

Version control is implemented using Git and GitHub. All project code, scripts, and configuration files are stored in a structured repository.

The repository includes:

\begin{itemize}
	\item Source code for data preprocessing and forecasting models
	\item Documentation describing the project setup and execution steps
	\item Configuration files specifying software dependencies
\end{itemize}

This documentation ensures that the project environment can be reproduced and that the experiments and results can be verified by other researchers.